\chapter{Filing an Issue}\label{cha:issues}\markboth{Filing an Issue}{Filing an Issue}

This machine has one author, and an issue you file here will, in all
likelihood, be handed straight to an AI coding session to work on. That
is worth knowing before you write one, because it changes what a good
report looks like.

So the form below is not really a form. It is a prompt. Filled in, it
carries enough for the work to start without a single question coming
back to you; left as prose, it usually costs one exchange to establish
things you already knew when you typed it.

\section{Three things first}
None of these is politeness. Each one removes a day.

\textbf{Type \texttt{DUMP} while it is wrong.} The machine writes its whole
state --- registers, the screen as you are looking at it, the last 4096
instructions, your last keystrokes, and the log of every command you have
typed --- to \pth{dumps/dump-NNN.txt} on the host. It is written for
exactly this purpose. If you can manage only one line of the report, make
it the dump: it is worth more than any description either of us could
write.

\textbf{Try it again with \texttt{k4510 {-}{-}no-startup.bat}.} A
\pth{/STARTUP.BAT} that has gone wrong looks exactly like a machine that
has gone wrong, and it has cost this project a diagnosis more than once.
Ten seconds of your time settles it.

\textbf{Look it up in this book.} If the machine disagrees with a page
here, that is still a fault worth filing --- it is simply this book's
fault rather than the machine's, and it gets fixed the same way. Say
which page, and the two halves of the project can sort out between them
which one was lying.

\section{The report}
Copy this, fill in what you can, delete what does not apply. Nothing in
it is required except the last two lines.

\begin{form}
Machine:  BMC-K4510, desktop  |  Raspberry Pi 3B+
Version:  what INFO -v gives, e.g. build 0.3-044bac5+  (a trailing
          + means that build was made from an edited tree).  Or,
          if it will not boot, the version on this book's cover
Where:    shell | EhBASIC | BBC BASIC | CP/M | Forth | Pascal |
          VI or EDIT | the F7 menu | this handbook

What I typed, exactly:

    LOAD "DEMO.BAS"
    RUN

What happened:
What I expected:
Why I expected it:  handbook p.NN | a real C64 or Beeb does it |
                    it worked before

Still wrong with  k4510 --no-startup.bat ?   yes | no | did not try
Dump:  dumps/dump-017.txt   (type DUMP, then attach that file)
\end{form}

\subsection*{Why those lines and not others}

\textbf{What I typed, exactly.} Not a summary of what you did --- the
lines themselves, as they went in. The machine is deterministic and its
figures in this book are captured by replaying keystrokes into it, so a
transcript can be run again as written. A paraphrase cannot.

\textbf{Why I expected it.} This is the line that most reports leave out
and most exchanges are then spent on. The K4510 is a fantasy machine
(page~\pageref{cha:disclaimer}): there is no silicon anywhere that it can
be measured against, so ``wrong'' only means anything next to what you
were expecting and where you got that from. It is also what separates a
bug from a design decision from an error in this book --- three different
repairs, and the answer decides which one starts.

\textbf{Where.} The machine is several machines: a shell, two BASICs, a
Z80 under CP/M, a Forth, a Pascal, two editors and a book. Naming yours
is what routes the report.

\section{Where it goes}
\begin{center}
\url{https://github.com/mlongval/bmc-k4510/issues}
\end{center}

\noindent The same block appears there as the repository's issue
template, generated from this page so that the two cannot drift apart. If
you have no account there, the same text in an email is worth exactly as
much.

Feature requests and disagreements are welcome in the same shape ---
``what I expected'' does the work whether or not anything is broken.

\begin{aside}
\textbf{One thing to check before you attach a dump.} It carries the
shell log, which is every command line you typed in that session, and
file names from your own disk along with them. It is a plain text file:
open it and have a look before it goes anywhere public.
\end{aside}
